\cotefile{157-5}
\scanpage{633}
\mspage{21}

à vérifier seulement

\[
 v^{*}v_{!}(TC)\subset TC,\tag*{(2.3.7)}
\]

or on a, pour $X=(X_i)_{i\in I_0}\in M(I_0)$

\[
\begin{aligned}
 v_!(X)(i)&\simeq\coprod_{j\in I_0}\unsure{v|u}_{j\,!}(X_j)(i)\\
 &\simeq\coprod_{j\in I_0}\underbrace{X_j^{(\operatorname{Hom}(i,j))}}_{\substack{\text{somme directe}\\\text{de }\operatorname{Hom}(i,j)\text{ copies}\\\text{de }X_j}},
\end{aligned}
\]

et la relation (2.3.7) résulte du fait que
$TC$, \inserted{\unsure{étant|il est} $Q$-clos à g.} est stable par petites sommes directes.

\uline{Prouvons $2^\circ$}. Les formules (2.3.2) sont
essentiellement tautologiques, les deux dernières
\unsure{résultent|résultant} de la troisième et de la deuxième,
en vertu de 1.5 appliqué au couple de
foncteurs adjoints $(u_!,u^*)$. La relation
$u_!(C\langle I\rangle)\subset C\langle I'\rangle$, et le fait que
$u_!$ transforme objet initial en objet initial,
impliquent (2.3.3). \unsure{Que} $u_!$ restreint à
$M(I)_c$ soit compatible aux localiseurs \unsure{provient}
de $u_!(TC\langle I\rangle)\subset TC\langle I'\rangle$, et du
fait que \unsure{tte|toute} flèche $f$ dans $M(I)_c$ qui est
dans $W(I)$, se décompose en $f=pi$, avec
$i\in TC$ et $p$ inverse à gauche d'une $j\in TC\langle I\rangle$.
Ceci \deleted{\unsure{dit}}\inserted{vu}, \deleted{\unsure{les}}\inserted{\unsure{l'une des|l'une}} deux premières formules
(2.3.2), montre qu'on est sous les conditions
% archive sheet number: 632
\scanpage{634}
\mspage{22}

du théorème d'adjonction de Quillen, d'où
l'adjonction de $u_!^W$ et de $u^*_{W}$, \unsure{cqfd|qed}.

\uline{Corollaire 2.4}. Soit \deleted{\unsure{$(M,W,C,F)$}} $W$ un localiseur
de $M$, tel qu'il existe un \deleted{\unsure{struct.}} triple
de Quillen spécial $(W,C,F)$ induisant $W$.
Alors le prédérivateur \deleted{défini} $\operatorname{Ho}_W$ \inserted{sur Cat} satisfait
aux axiomes Der 1 (décomposition), Der 2
(localisation), Der 3 (objet final), Der 4 bis
(\unsure{existence|existences} des $u_!^W$), Der 5 bis (calcul des
fibres des $u_!^W$), Der 6 bis (exactitude
à gauche, \deleted{\unsure{?}}\inserted{i.e. exactitude à g.} des carrés \inserted{$W$-}cocartésiens \unsure{dans|des} les $\operatorname{Ho}_W(1)$),
Der 6 (exactitude à droite \deleted{, des} carrés $W$-cartésiens).

On vient de prouver Der 4 bis, et Der. 5 bis
est \inserted{pratiquement} contenu dans le mode d'obtention de \deleted{\unsure{cet}}
$u_!$, compte tenu de la prop. 1.12, qui implique

\deleted{\uline{Corollaire \unsure{2.5}}. Soit $w\colon J\unsure{-|\longrightarrow} I$ une flèche de
Cat, supposons que l'on soit dans l'un des
trois cas suivants

a) $w$ est \unsure{coétale}

b) $J$ est discrète

c) $\forall\ i\in I$, $\deleted{\unsure{C}}J_{/i}$ est \unsure{bisomme|somme|réunion} de catégories ayant chacune un objet final.}
que pour \deleted{\unsure{$i'\text{ de }I'$}}
\deleted{\illegible{1}} $X$ dans $M(I)_c$, et $i'$ dans $I'$, l'image inverse
de $X$ dans $M({}_{\unsure{i'|i}\backslash}I)$ est dans $M({}_{i'\backslash}I)_c$ \ldots.
% archive sheet number: 633
\scanpage{635}
\mspage{23}

L'axiome de décomposition est satisfait pour
tout prédérivateur de la forme $\unsure{\mathbb{D}|\underline{D}}_{M,W}$, \deleted{\illegible{1}}
\deleted{\unsure{$W$}} l'axiome d'exactitude à gauche \inserted{Der 5 bis} est valable,
pourvu que les $(M(I),W(I))$ \unsure{s'inscrivent|s'insèrent} dans
une structure à cofibrations, ce qui est le cas.
De $=$ pour l'axiome Der 2 de localisation,
pourvu que de plus $W$ soit fortement
saturé, ce qui est également le cas. Enfin,
l'axiome de l'objet final dit que si
$I$ admet un objet final, alors

\[
 \operatorname{Ho}_W(e)\unsure{\;-\;|\longrightarrow}\operatorname{Ho}_W(I)
\]

est pl. fidèle, ou encore que le morphisme
d'adjonction

\[
 u_!u^*(\unsure{\xi|\zeta})\longrightarrow \xi
\]

est iso. Mais en vertu de \unsure{1.12}, cas $3^\circ$,
dans ce cas $u^*$ \deleted{\unsure{transforme}}\inserted{applique} $M(e)_c\simeq M_c$ dans
$M(I)_c$, donc le calcul du couple de
foncteurs adjoints $u_!^W,u^*_W$ peut se faire en
utilisant les \unsure{restrictions} de $u_!,u^*$ à $M(I)_c$ et
$\unsure{M_c}$ respectivement. Comme le morphisme
d'adjonction $u_!^c u_c^*\unsure{\longrightarrow|\mapsto}\mathrm{id}$ est iso, il en
est donc de $=$ de $u_!^W u^*_W$ qui s'en
déduit en localisant, \unsure{qed|cqfd}.

2.5. \unsure{Pour} achever de construire $\operatorname{Ho}_W$ comme
un dérivateur sur Cat \deleted{tout} entier, il faut
% archive sheet number: 634
\scanpage{636}
\mspage{24}

\unsure{donc} établir Der 4, Der 5 : existence des $u_*^W$, et
calcul des fibres de $u_*^W(X)$. Quant à Der 6,
il est acquis automatiquement, du fait
que chaque $(M(I),W(I))$ s'insère dans une
structure de fibrations $(W(I),F(I))$. (Oublié
dans la démonstration de 2.4.) L'idée naturelle
est d'utiliser à présent la structure de
triple

\[
 \left(W(I),C(I),\underset{\substack{\|\,\mathrm{def}\\ TC(I)_{*}}}{F[I]}\right),\tag*{(2.5.1)}
\]

mais la difficulté est de prouver la \unsure{factorisabilité|factorisabili-}
pour les couples

\[
 C(I)\underset{\unsure{<|?}}{\leftrightarrow}TF[I],\quad
 \underset{\substack{\|\,\mathrm{def}\\ C(I)\cap W(I)}}{TC(I)}\longleftrightarrow F[I],\tag*{(2.5.2)}
\]

faute de pouvoir montrer, dans les conditions
générales de 2.3 seulement (triple de Quillen
spécial) que $C(I)$ et $TC(I)$ sont également
spéciaux. (Cf. \unsure{question} 1.11.) Mais
moyennant des conditions convenables sur
$M$, on arrive à s'en tirer si $C=\underline{\mathrm{mono}}$.
L'énoncé pertinent, un peu technique, se
trouve dans \uline{XV} cor 7.4 (p 98--99), que je
vais rappeler ici.
% archive sheet number: 635
\scanpage{637}
\mspage{\unsure{25}}

\uline{Théorème} 2.6. Soit $M$ une catégorie satisfaisant les
conditions suivantes

(2.6.1) \(\Bigg\{\)
a) $M$ est spéciale, \unsure{i.e.} stable par petites $\varinjlim$,
et ses éléments sont accessibles.

% slanted
\margin{\uline{NB} a $+$ c signifient que $M$ est une ``catégorie $U$-paratopos'', et impliquent l'existence d'une filtration cardinale satisfaisant à d) (sauf \unsure{l'assertion|l'accessibilité} \unsure{concernant les quotients} $Z'$ de $Z$) \unsure{cf.|c'est} \uline{\unsure{cat. de faisceaux !|de bien voir !}} Donc il s'agit d'un $U$-paratopos \unsure{satisfaisant en plus aux conditions} b) e) f) [\unsure{Est-on alors|Est-ce alors} \unsure{un hypo ?|une hypo ?} \uuline{Non}, \unsure{c'est|car} e), f)].}

b) \deleted{les limites \unsure{inductives}} filtrantes dans $M$ sont exactes à gauche \unsure{(}et dans $M$ les \deleted{\unsure{lim}} $\varprojlim$ finies sont représentables \unsure{(résulte de|résultent de} $a+c$)

c) $M$ admet une filtration cardinale
$(\mathrm{Filt}^{\pi})_{\pi\geq\pi_0}$.

d) les $\mathrm{Filt}^{\pi}$ sont stables par $\varinjlim$ indexées
par des cat. $I$ telles que $\operatorname{card}\mathrm{Fl}(I)\leq\pi$
et par quotients ($Z\in\mathrm{Filt}^{\pi}$ et $Z\longrightarrow Z'$
épi implique $Z'\in\mathrm{Filt}^{\pi}$)

\margin{\uline{NB} La condition d) est \unsure{sans doute inutile}, du moins si les épi sont effectifs, car d) sera satisfaite en tous cas pour $\pi=2^{\unsure{c|\ell}}$, avec $c\geq\pi_0$}

e) Si $A_1,A_2$ sont deux sous-objets d'un
objet $B$ de $M$, et $A_0=A_1\cap A_2$, alors
$A_1\amalg_{A_0}A_2\longrightarrow B$ est un mono.

% slanted
\margin{\uline{NB} La condition e) n'est \uline{pas} satisfaite pour la catégorie (Gr), et le th. est en défaut pour cette catégorie, laquelle n'admet \uline{pas} ``assez d'injectifs'' cf \deleted{\unsure{S}} SGA 4 I 8.17.9. Par contre, e) et f) sont satisfaites pour les catégories \uline{abéliennes}.}

f) Si on a un carré cocartésien

\diagram{\begin{tikzcd}[ampersand replacement=\&,row sep=small]
A_0 \arrow[r,"i"] \arrow[d] \& A_1 \arrow[d] \\
A_2 \arrow[r,"i'",no head] \& A_3
\end{tikzcd}}

avec $i$ un mono, alors $i'$ est mono.

% slanted
\margin{\uline{XVIII} cf Th. 6.7 (p. 157) pour une formulation plus \unsure{symétrique|synthétique} de l'ensemble des conditions a) à f) pour $M$.}

(On dira alors que $M$ est ``\uline{très spéciale}''.) Soit
de plus

\[
 W\subset\mathrm{Fl}\,M\tag*{(\unsure{2.6.2})}
\]

un localiseur de $M$, qui satisfait les conditions

(\unsure{2.6.3}) \(\Bigg\{\)
a) $W$ est stable par $\varinjlim$ filtrantes

% slanted
\margin{\uline{NB} a) implique déjà que $W=W^{\unsure{\natural|\sharp}}$}

b) $W$ est de plus hypo\uline{accessible}, i.e. il
existe une \uline{\unsure{petite}} partie $W_0\subset W$,
telle que $W$ soit la hypoenveloppe de
\unsure{$W$|$W_0$}, \inserted{pour $\unsure{\mathscr{V}|U}$ \unsure{sig.}} \unsure{i.e.} la plus petite partie de
$\mathrm{Fl}\,M$ \deleted{stable} contenant $W_0$, et stable
par $\varinjlim$ filtrantes.
% archive sheet number: 636
\scanpage{638}
\mspage{26}

Posons alors

\[
 C=\operatorname{Mono}(M),\quad TC=C\cap W,\quad TF=C_*,\quad F=(TC)_*.\tag*{(2.6.4)}
\]

Sous ces conditions

$1^\circ$) $C\mathrel{\leftarrow\!\!|\!\!\rightarrow}TF$, plus précisément $C\unsure{\mathrel{\boldsymbol{\leftarrow\!\!|\!\!\rightarrow}}|\mathrel{\Leftarrow\!\!|\!\!\Rightarrow}}TF$
est un couple de Quillen spécial.

$2^\circ$) Pour que $TC\longleftrightarrow F$ soit un couple
de Quillen, il f. et s. que $TC$ soit
stable par cochangement de base, et alors
c'est un couple de Quillen spécial.

$3^\circ$) Pour que $(W,C,F)$ soit un triple de
Quillen (lequel sera nécessairement spécial
par $1^\circ,2^\circ$, \deleted{\unsure{donc}} et en particulier \unsure{clos}) il
f. et s. que les conditions suivantes soient
satisfaites

\[
\left\{\begin{aligned}
 \text{a)}\quad & TC\bigl(\overset{\mathrm{def}}{=}W\cap\mathrm{Mono}\bigr)\ \text{est stable par cochange-}\\
 &\quad\text{ment de base}\\
 \text{b)}\quad & TF\bigl(\overset{\mathrm{def}}{=}(\mathrm{Mono})_*
 \unsure{=|\underset{\vert}{=}}\text{ ``applications de}\\
 &\hspace{8em}\text{type ``injectif''}\bigr)\subset W.
\end{aligned}\right.\tag*{(2.6.5)}
\]

% slanted
\margin{\uline{NB} les conditions (2.6.5) équivalent à dire que $(W,C)$ est une structure de \inserted{\unsure{cat. à}} cofibrations sur $M$.}

et la condition b) équivaut encore à \deleted{celle}.

(2.6.6) \uline{b'}) Toute flèche $f$ de $M$ (ou seulement
de $TF$) se factorise en $f=pi$, avec
$i$ mono, $p\in W$. \deleted{\unsure{(On verra en}}

\uline{Corollaire 2.7.} \deleted{\illegible{1}}\inserted{$M$ et $W$ étant comme dans (2.6.1)\deleted{\unsure{a}}, (2.6.3),}
\deleted{\unsure{présence de}} Pour que $(W,C,F)$ soit un triple de
Quillen \uline{propre à gauche} (i.e. $C\subset\operatorname{Cof}_W$), il
f. et s. que $W$ satisfasse la condition
(2.6.5\deleted{\unsure{a}} b) (i.e. \deleted{\unsure{$TF$}} $W$ contient les \unsure{flèches} de
type injectif) et de plus
% archive sheet number: 637
\scanpage{639}
\mspage{27}

(\unsure{2.7.1}) $\deleted{\unsure{\mathrm{Mono}}}C\subset\operatorname{Cof}_W$, i.e. \deleted{pour \unsure{tout|tt} diag}\inserted{la condition}

\unsure{a'}) Pour tt diagramme cocartésien

\diagram{\begin{tikzcd}[ampersand replacement=\&,row sep=small]
A_0 \arrow[r,"i",no head] \arrow[d,"f"] \& A_1 \arrow[d,"f'"]\\
A_2 \arrow[r,"i'",no head] \& A_3
\end{tikzcd}}

avec $i$ mono, si $f\in W$, alors $f'\in W$.
(Stabilité de $W$ par \unsure{cochangemt} de base
par monomorphismes.)

Sous ces conditions, ce triple de Quillen est
spécial, et a fortiori clos.

Ce corollaire résulte de 2.6. $3^\circ$), puisque
$C\subset\operatorname{Cof}_W$ implique $TC\overset{\mathrm{def}}{=}C\cap W=C\cap W^{\unsure{\mathrm{univ}|\mathrm{uni}}}$,
\deleted{donc que} (car $\operatorname{Cof}_W\cap W=W^{\mathrm{univ}}$), donc que
$TC$ est stable par \unsure{coch.} de base, d'où la
condition (2.6.5 a).

2.8. On dit qu'une \uline{structure de Quillen}
$(W,C,F)$ sur $M$ est \uline{très} \uline{spéciale}, si elle
est obtenue comme au th. 2.1, i.e.
si on a les conditions

\[
\left\{\begin{aligned}
\text{a)}\quad & M\ \text{est très spéciale, i.e. satisfait aux}\\
 &\text{conditions (2.6.1)}\\
\text{b)}\quad & W\ \text{est stable par }\inserted{\text{petites}}\ \varinjlim\ \text{filtrantes et}\\
 &\text{est accessible (conditions (2.6.3))}\\
\text{c)}\quad & C=\mathrm{Mono}.
\end{aligned}\right.\tag*{(2.8.1)}
\]

Donc pour qu'un localiseur $W$ \unsure{de|dans} $M$ corresponde
à une structure de Quillen \deleted{spé} très spéciale
% archive sheet number: 638
\scanpage{640}
\mspage{28}

dans $M$ (laquelle sera forcément unique), il f. et
s. que $M$ et $W$ satisfassent les conditions
(2.8.1 a) b)), et que de plus on ait

\[
\left\{\begin{aligned}
\text{a)}\quad &TC\overset{\mathrm{def}}{=}\mathrm{Mono}\cap W\ \text{est stable par cochg}^{\mathrm t}\text{ de base}\\
\text{b)}\quad &TF\overset{\mathrm{def}}{=}(\mathrm{Mono})_*\subset W,\quad\text{i.e. }W\text{ contient}\\
 &\text{les flèches de type injectif dans}\\
 &M.
\end{aligned}\right.\tag*{(2.8.2)}
\]
\deleted{\unsure{(2.8.3)}}

% slanted
\margin{On dit \unsure{aussi|ainsi} que le localiseur $W$ est très spécial, s'il définit une str. de Quillen très spéciale.}

On trouve maintenant le

\uline{\deleted{\unsure{Corollaire}}Théorème} 2.9. Soit $(W,C\unsure{=|\underset{\vert}{=}}\mathrm{Mono},F=(TC)_*=(W\cap\mathrm{Mono})_*)$
un triple de Quillen très spécial dans une
catégorie $M$ (2.8.). Alors \deleted{pour}

$1^\circ$) Pour \unsure{tout} $I$ dans Cat, $(W(I),\deleted{\unsure{W}}C(I),F[I])$
est un \deleted{couple}\inserted{triple} de Quillen très spécial
dans $M(I)$, et on a \deleted{\unsure{si $v\colon I_0\unsure{-|\longrightarrow}I$ est} défini \unsure{comme}}\inserted{en particulier}

\[
 C(I)=\operatorname{Mono}(M(I)),\tag*{(2.9.1)}
\]

et de plus

\[
 F[I]\cap W(I)=(TF)[I].\tag*{(2.9.2)}
\]

$2^\circ$) Pour toute flèche $u\colon I\unsure{-|\longrightarrow}I'$ dans Cat,
et \deleted{\unsure{où|au}} \deleted{\unsure{le}}le couple de foncteurs adjoints

\diagram{\begin{tikzcd}[ampersand replacement=\&]
\text{(2.9.3)} \& M(I) \arrow[r,shift right=0.8ex,"u_*"'] \& M(I') \arrow[l,shift right=0.8ex,"u^*"'] ,
\end{tikzcd}}

on a

\[
\left\{\begin{aligned}
 u^*(W(I'))&\subset W(I)\\
 u^*(C(I'))&\subset C(I)\\
 u^*(TC(I'))&\subset TC(I)\\
 u_*(F[I])&\subset F[I']\\
 u_*(TF[I])&\subset TF[I']
\end{aligned}\right.,\tag*{(2.9.\unsure{4|3})}
\]

enfin $u_*$ induit un foncteur
% archive sheet number: 639
\scanpage{641}
\mspage{29}

\[
 u_*^f\colon M(I)_f\longrightarrow M(I')_f\tag*{(2.9.5)}
\]

\margin{sous-catégories pleines des objets fibrants}

\uline{compatible aux localiseurs} $W(I)_f$ et $W(I')_f$
induits par $W(I),W(I')$ respectivement, donc
induit un foncteur

\[
 u_*^W\colon\operatorname{Ho}_W(I)\unsure{-|\longrightarrow}\operatorname{Ho}_W(I').\tag*{(2.9.6)}
\]

Ce foncteur est un \uline{adjoint à droite} du
foncteur \deleted{\unsure{$u_W^*$}}

\[
 u^*_W\colon\operatorname{Ho}_W(I')\mathrel{-}\operatorname{Ho}_W(I)\tag*{(2.9.7.)}
\]

induit par $u^*$ (2.9.3).

\uline{Démonstration de} 2.9. Pour prouver $1^\circ$),
il faut vérifier \inserted{d'abord} que les conditions préliminaires
(2.8.1 a) b)) sur $M,W$ sont stables
par passage de $M,W$ à $M(I),W(I)$, ce
qui se fait par \unsure{AQT}. \deleted{De plus, \unsure{et de} \unsure{Au théorème 2.6|Au théorème 2.7}
présence des conditions (2.8.2 a) b)), le \unsure{fait}
théorème 2.6 implique alors que (pour \unsure{le|v}}
\deleted{\unsure{Mono}} D'ailleurs on vérifie tout de suite que
\deleted{\unsure{$W(I)$}} $C(I)=\operatorname{Mono}(M(I))$,
et le th. 2.6 ($1^\circ,2^\circ$) implique \inserted{alors} que

\[
\begin{aligned}
 &\deleted{(C(I),C(I)_*}\\
 C(I)&\longleftrightarrow C(I)_*\overset{\mathrm{def}}{=}TF[I]\\
 \underset{\substack{\|\,\mathrm{def}\\ C(I)\cap W(I)}}{TC(I)}
 &\longleftrightarrow(TC(I))_*\overset{\mathrm{def}}{=}F[I]
\end{aligned}
\]

sont des couples de Quillen spéciaux, \unsure{(Étant|étant}
entendu, pour le deuxième, que $TC(I)$ est bien
% archive sheet number: 640
\scanpage{642}
\mspage{30}

stable par cochangement de base, $TC$\deleted{\unsure{f}} l'étant.) Pour
vérifier que l'on a un triple de Quillen
très spécial, il \deleted{suffit de}\inserted{reste, en vertu de 2.8, à} vérifier que

\[
 TF[I]\subset W(I)
\]

Ce qui \deleted{est} s'écrit

\[
 v^*(TF[I])\subset W(I_0),
\]

où

\[
 v\colon I_0\unsure{-|\longrightarrow}I
\]

est défini comme dans (1.12) (p. \unsure{10|20}). En fait,
on aura même

\[
 v^*(TF[I])=v^*v_*(TF(I_0))\subset
 \underset{\substack{\|\\ F(I_0)\cap W(I_0)}}{TF(I_0)},
\]

et la démonstration est duale de celle de

\[
 v^*v_!(TC(I_0))\unsure{\subset|\in} TC(I_0),
\]

(cf (2.3.\unsure{7|2}) p. 21), elle ne dépend pas des
questions de factorisation dans $M(I)$, et provient
ici de la stabilité de $TF$ par petits
produits. Cela prouve $1^\circ$.

\deleted{Une fois}
La démonstration de $2^\circ$) est maintenant
\inserted{la duale de} celle de (2.3. $2^\circ$). Comme cela nous fait
\deleted{\unsure{prouver}}\inserted{répéter} deux fois la $=$ démonstration, cela
montre qu'on n'a pas ``monté la \unsure{suite|sauce|boîte}'' \unsure{tout}
à fait \unsure{avec} assez de soin. Donc à titre
de \unsure{rattrappage|rattrapage}, j'ajoute un
% archive sheet number: 641
\scanpage{643}
\mspage{31}

\uline{Scholie 2.10} Soit $(W,C,F)$ un triple de Quillen
clos dans $M$. Pour \deleted{\unsure{un}} $I$ dans Cat, pour que
$(W(I),C\langle I\rangle,F(I))$ soit un triple de Quillen
clos pour $\unsure{M(I)}$, il f et s. que les
couples $C\langle I\rangle\longleftrightarrow TF(I)$ et $TC\langle I\rangle\longleftrightarrow F(I)$
soient factorisants. Supposons que cette
condition soit satisfaite pour $I$ et pour $I'$,
et soit $u\colon I\longrightarrow I'$ une flèche \deleted{\unsure{dans ?}}\inserted{de}
$I$ dans $I'$. Alors les assertions de 2.3. $2^\circ$
sont valables, en particulier on trouve un
foncteur $u_!^W\colon\operatorname{Ho}_W(I)\unsure{-|\longrightarrow}\operatorname{Ho}_W(I')$, adjoint
à gauche de $u^*_W$.

% slanted
\margin{Dire aussi que le calcul des fibres de $u_!^W(X)$, prévu dans Der 5 bis, est OK.}

Ces résultats se
dualisent aussitôt (par $M\unsure{\mapsto|-} M^{\unsure{o|\circ}}, I\unsure{-|\mapsto}I^{o}$)
pour les triples $(W(I),C(I),F[I])$,
et pour la construction de $u_*^W$, adjoint
à droite de $u^*_W$. En particulier, pour que
le triple précédent soit un triple de Quillen
(forcément clos) il f. et s. que les deux
couples $C(I)\longleftrightarrow TF[I]$ et $TC(I)\longleftrightarrow F[I]$
soient factorisants ; et quand cette condition
est satisfaite pour $I$ et pour $I'$, \deleted{\unsure{et}}(pour une
flèche $u\colon I\mathrel{-} I'$ dans Cat donnée), alors
$u_*^W$ existe et se construit dualement de
la construction de $u_!^W$ explicitée dans
2.3. $2^\circ$.
% archive sheet number: 642
\scanpage{644}
\mspage{32}

\uline{Corollaire} 2.11. Le prédérivateur $\unsure{\mathbb{D}|\underline{D}}_{M,W}$ défini
par un localiseur \uline{très spécial} $W$ dans $M$
est un dérivateur au plein sens du terme,
i.e. il satisfait à tous les axiomes
Der 1 à Der 6 (y inclus Der 4 bis et
Der 5 bis et Der 6 bis) des dérivateurs.

\uline{NB} À ma propre surprise, je n'ai pas été capable
de \unsure{trouver|donner} un cas d'un localiseur $W$,
qui définisse un dérivateur sur Cat tout
entier (avec existence des $u_!^W,u_*^W$),
\uline{que} : pour $W$ un localiseur très spécial,
donc \unsure{s'inscrivant|s'insérant} dans \unsure{:} un triple de
Quillen très spécial -- ou quand on
peut s'y ramener de la façon \unsure{standart|standard}.

(Par exemple dans le cas \deleted{\illegible{1}}\inserted{où} $M=\mathrm{Cat}$, $W$
localiseur fondamental \inserted{satisfaisant Loc \unsure{5|6} a) et Loc 8, 8 bis} -- les
mono $\in W$ se refusent formellement à
être stables par cochangement de base !)

% slanted
\margin{axiome des carrés cocartésiens (version directe), axiome des limites inductives, accessibilité}

J'ignore totalement si le résultat mirabolant
``prouvé'' par Anderson est vrai ou faux. \unsure{Ça|Sa} a l'air
\inserted{presque} trop beau pour être vrai \ldots. D'ailleurs, je ne
connais pas de structure de Quillen spéciale qui
ne soit très spéciale -- désolé !
% archive sheet number: 643
\scanpage{645}
\mspage{33}

3. \uline{Cas des catégories à cofibrations}, \uline{des cat. à fibrations}.

Je rappelle les axiomes d'une \uline{structure à cofibrations}
$(W,C)$ sur une catégorie $M$.

$\mathrm{Cof}_0$ : $W$ est un localiseur.

$\mathrm{Cof}_1$ : $C$ contient les iso, est stable par \unsure{cochangemt}
de base, par composition.

$\mathrm{Cof}_2$ : $C\cap W$ est stable par \unsure{cochangement} de base
(et aussi par composition, et contient les
iso, en vertu de $\mathrm{Cof}_0,\mathrm{Cof}_1$)

$\mathrm{Cof}_3$ : Toute flèche $f$ \unsure{dans|de} $M$ se factorise en
$f=pi$, avec $i\in C,p\in W$.

(Axiomes de K. S. Brown, légèrement généralisés
par Anderson.) Il y a un ``axiome de
continuité'' \inserted{ou ``\uline{\unsure{axiome des fins|axiome des feux|axiome sans fin}}''} facultatif, qui est le suivant

$\mathrm{Cof}_4$ : \inserted{$M$ \unsure{stable} par petites $\varinjlim$ filtrantes, et}
$C$ et $TC\overset{\mathrm{def}}{=}C\cap W$ sont stables par
composition transfinie. \deleted{\unsure{ou}}.

Ainsi, $\mathrm{Cof}_1$, $\mathrm{Cof}_2$, $\mathrm{Cof}_4$ peuvent se résumer
en disant que $C$ et $TC$ sont $Q$-saturés
\unsure{transfiniment} à gauche.

La structure de \unsure{cofibration|cofibrations} est dite
\uline{spéciale} si elle, \deleted{est stable} satisfait $\mathrm{Cof}_4$,
si $M$ est spéciale, et si $C$ admet
une \uline{petite} partie $C_0$ qui soit $Q$-génératrice
à gauche, i.e. telle que

\[
 C\subset \unsure{C_0^{\sim}|\widetilde{C_0}}\quad\text{i.e.}\quad\widetilde C=C_0^{\sim}.\tag*{(3.1)}
\]
% archive sheet number: 644
\scanpage{646}
\mspage{3\deleted{\unsure{5}}\unsure{4|5}}

Alors, par le th. de factorisation 1.1, on a

\[
 C\mathrel{\leftarrow\!\!|\!\!\rightarrow}TF\overset{\mathrm{def}}{=}C_*,
\]

et l'axiome de factorisation $\mathrm{Cof}_3$ implique
(cf 1 \unsure{M|11}) $TF\subset W^{\unsure{\natural|\sharp}}$. Plus généralement\unsure{.}

\uline{Prop. 3.1} Supposons \inserted{que $M$ soit \unsure{spécial|spéciale} (stable par petites $\varinjlim$, et éléments accessibles),} que \deleted{\unsure{f}} $(W,C)$ satisfasse
à $\mathrm{Cof}_0$, $\mathrm{Cof}_1$, et que $C$ soit stable par
composition transfinie (ce qui, joint à $\mathrm{Cof}_1$,
\unsure{s'énonce|se résume} en disant que $C$ est $Q$-saturé
transfini à gauche), enfin que $C$
admette une petite partie $C_0$, $Q$-génératrice
à gauche, de sorte que $C\mathrel{\leftarrow\!\!|\!\!\rightarrow}TF\overset{\mathrm{def}}{=}C_*$.

\deleted{Supposons qu'on ait $W=W^{\natural}$, i.e. $W$ stable par
facteurs directs}. Alors, \inserted{pour que} l'axiome de
factorisation $\mathrm{Cof}_3$ \deleted{équivaut à $TF\subset W$}\inserted{soit satisfait, il suffit qu'on ait}

\[
 TF\subset W,\tag*{(3.1.1)}
\]

et il faut qu'on ait $TF\subset W^{\natural}$ ; donc si $W=W^{\natural}$,
l'axiome $\mathrm{Cof}_3$ équivaut à (3.1.1).

La structure de cofibrations est dite \inserted{\uline{semi}-}\uline{close}
si \inserted{$C$ est $Q$-clos à g., i.e. si} on a $C=\unsure{C^{\sim}|\widetilde C}$, (ce qui implique
que $C$ est stable par composition transfinie),
\uline{close} si de plus $TC=(TC)^{\sim}$. (ce qui
implique l'axiome transfini $\mathrm{Cof}_4$).
% archive sheet number: 645
\scanpage{647}
\mspage{35}

Je ne reviens pas ici sur la théorie de K. S. Brown
des catégories à cofibrations. C'est à ma connaissance
l'ensemble d'axiomes le plus économique
connu, pour assurer d'une part une bonne
prise sur $\operatorname{Ho}_W$ (par une notion convenable
\unsure{d'|de'}homotopie d'applications, et un calcul \unsure{de|des}
fractions dans $\pi M_c$), et d'autre part
\deleted{la}\inserted{l'existence de la} suite \uline{exacte} \unsure{de cofibrations|de cofibration|des cofibrations}, donc une
théorie des cofibres $W$-homotopiques. [NB
Quand la structure est \uline{propre}, i.e. $C\subset\operatorname{Cof}_W$,
alors l'existence de cette structure revient
à dire simplement que ``$W$ admet assez de
$W$-cofibrations'', i.e. que toute $f\in\mathrm{Fl}\,M$
se factorise en $pi$, avec $i\in\operatorname{Cof}_W,p\in W$ ;
et on peut alors prendre $C=\operatorname{Cof}_W$. Dans le
cas où $W$ est stable par \inserted{petites} $\varinjlim$ filtrantes,
$\unsure{C|C_*}$\inserted{$=\operatorname{Cof}_W$}, et $TC=W^{\unsure{\mathrm{univ}|\mathrm{uni}}}$ satisfont \inserted{de plus} l'axiome transfini
$\mathrm{Cof}_4$. \deleted{\unsure{Ces}}] Enfin, K. S. Brown montre
que toute flèche $i\colon X\longrightarrow Y$ dans $C$,
\uline{avec} $X,Y$ \uline{cofibrants} (i.e $\varnothing_M\unsure{-|\longrightarrow}X,\varnothing_M\unsure{\longrightarrow|-}Y$
dans $C$) est dans $\operatorname{Cof}_{W^{\unsure{\natural|\sharp}}}$ (``lemme de
Brown'').

Je m'intéresse ici à la question
d'étendre la structure à cofibrations
à une catégorie de la forme $M(I)$,
et d'utiliser cette extension pour construire
% archive sheet number: 646
